% 半人马小行星
% license CCBYSA3
% type Wiki

本文根据 CC-BY-SA 协议转载翻译自维基百科\href{https://en.wikipedia.org/wiki/Centaur_(small_Solar_System_body)}{相关文章}。

\begin{figure}[ht]
\centering
\includegraphics[width=6cm]{./figures/4207f1463529dc38.png}
\caption{} \label{fig_Centau_2}
\end{figure}
在行星天文学中，半人马天体（centaur）是指一类在木星与海王星之间绕太阳运行、并且其轨道会与一颗或多颗巨行星轨道相交的小型太阳系天体。由于这一特性，半人马天体的轨道通常是不稳定的；几乎所有半人马天体的轨道动力学寿命都只有数百万年$^{[1]}$，但有一个已知的例外，即 514107 Kaʻepaokaʻāwela，其轨道可能是稳定的（尽管是逆行轨道）$^{[2]}$（见注释$^{[1]}$）。半人马天体通常同时表现出小行星和彗星的特征。它们的名称来源于神话中的半人马，即兼具马与人形态的生物。由于观测上对大尺寸天体存在偏向性，半人马天体的总体数量难以确定。对直径大于 1 km 的半人马天体数量估计范围从约 44,000 个$^{[1]}$到超过 1,000 万个$^{[4][5]}$不等。

按照喷气推进实验室（JPL）的定义（也是本文采用的定义），第一个被发现的半人马天体是 944 Hidalgo，发现于 1920 年。然而，直到 1977 年发现 2060 Chiron 之后，半人马天体才被视为一个独立的天体族群。目前已确认的最大半人马天体是 10199 Chariklo，其直径约为 250 km，大小相当于一颗中等尺寸的主带小行星，并且已知其拥有一套环系统。该天体发现于 1997 年。

目前尚无任何半人马天体被近距离拍摄过，不过有证据表明，土星的卫星菲比（Phoebe）——由卡西尼号探测器于 2004 年拍摄——可能是一颗起源于柯伊伯带并被俘获的半人马天体$^{[6]}$。此外，哈勃空间望远镜也获取了一些关于 8405 Asbolus 表面特征的信息。

谷神星可能起源于外行星区域$^{[7]}$，如果这一假设成立，它可以被视为一颗前半人马天体，但目前观测到的半人马天体均起源于其他区域。

在已知占据半人马式轨道的天体中，约有 30 个显示出类似彗星的尘埃彗发，其中 2060 Chiron、60558 Echeclus 和 29P/Schwassmann–Wachmann 1 在完全位于木星轨道之外的情况下仍表现出可探测的挥发物释放$^{[8]}$。因此，Chiron 和 Echeclus 同时被归类为半人马天体和彗星，而 Schwassmann–Wachmann 1 一直被归类为彗星。其他半人马天体（如 52872 Okyrhoe）也被怀疑曾出现过彗发。理论上，任何被扰动并足够接近太阳的半人马天体，最终都可能演化为彗星。
\subsection{分类}
若一颗天体的近日点距离或半长轴位于外行星之间（即木星与海王星之间），即可被归类为半人马天体。由于该区域内轨道在长期尺度上的固有不稳定性，即便是当前并未与任何行星轨道相交的半人马天体（如 2000 GM137 和 2001 XZ255），其轨道也会逐渐演化，并最终受到扰动，从而开始与一颗或多颗巨行星的轨道相交$^{[1]}$。一些天文学家仅将半长轴位于外行星区域的天体归为半人马天体；而另一些学者则认为，只要近日点位于该区域内即可归类为半人马天体，因为这两类轨道在动力学上同样是不稳定的。
\subsubsection{判据差异}
然而，不同机构在对边界天体进行分类时采用了不同的判据，这些判据通常基于轨道要素的具体取值：
\begin{itemize}
\item 小行星中心（Minor Planet Center，MPC）将半人马天体定义为近日点位于木星轨道之外（5.2 AU < q），且半长轴小于海王星半长轴（a < 30.1 AU）的天体$^{[9]}$。MPC 有时会将半人马天体与散射盘天体列为同一类群$^{[10]}$。
\item 喷气推进实验室（Jet Propulsion Laboratory，JPL）采用了相近的定义，将半人马天体界定为半长轴 a 介于木星与海王星之间的天体（5.5 AU ≤ a ≤ 30.1 AU）$^{[11]}$。
\item 相比之下，深黄道巡天（Deep Ecliptic Survey，DES）使用的是一种动力学分类方案。该分类基于将当前轨道向未来延拓 1000 万年时，其轨道行为在数值模拟中的变化。DES 将半人马天体定义为非共振天体，其瞬时（密切）近日点在模拟过程中任意时刻小于海王星的瞬时半长轴。该定义旨在与“行星交叉轨道”等价，并暗示这类天体在当前轨道上的寿命相对较短$^{[12]}$。
\item 《海王星之外的太阳系》（The Solar System Beyond Neptune，2008）一书将半长轴位于木星与海王星之间、且相对于木星的 Tisserand 参数大于 3.05 的天体定义为半人马天体；将相对于木星的 Tisserand 参数低于该阈值、并且为排除柯伊伯带天体而额外施加一个任意的近日点限制（q ≤ 7.35 AU，位于通往土星距离的一半处）的天体归类为木星族彗星；而将半长轴大于海王星、且处于不稳定轨道上的天体归类为散射盘天体$^{[13]}$。
\item 还有一些天文学家倾向于将半人马天体定义为：不与行星处于共振状态、近日点位于海王星轨道以内、并且可以证明在未来 1000 万年内很可能进入某一气体巨行星希尔球范围的天体$^{[14]}$。在这种定义下，半人马天体可被视为向内散射的天体，其与行星的相互作用更强、轨道演化也比典型的散射盘天体更快。
\item JPL 小天体数据库目前列出了 910 颗半人马天体$^{[15]}$。此外，还有 223 个跨海王星天体（即半长轴大于海王星，30.1 AU ≤ a），其近日点距离小于天王星轨道（q ≤ 19.2 AU）$^{[16]}$。
\end{itemize}
\subsubsection{歧义天体}
Gladman 与 Marsden（2008）$^{[13]}$ 的判据会将一些天体归类为木星族彗星：Echeclus（q = 5.8 AU，$T_J = 3.03$）和 Okyrhoe（q = 5.8 AU，$T_J = 2.95$）在传统上一直被归类为半人马天体。Hidalgo（q = 1.95 AU，$T_J = 2.07$）传统上被视为一颗小行星，但在 JPL 的分类中被列为半人马天体，在该判据下也会改变类别而被归为木星族彗星。29P/Schwassmann–Wachmann（q = 5.72 AU，$T_J = 2.99$）则根据所采用的定义不同，被同时归类为半人马天体或木星族彗星$^{[]}$。

其他处于不同分类方法分歧中的天体还包括 (44594) 1999 OX3，其半长轴为 32 AU，但会穿越天王星和海王星的轨道。在深黄道巡天（DES）中，它被列为外侧半人马天体。在内侧半人马天体中，(434620) 2005 VD 的近日点距离非常接近木星，其在 JPL 和 DES 中均被列为半人马天体。

一项近期关于柯伊伯带天体经由半人马区域演化的轨道模拟研究$^{[4]}$ 识别出一个位于 5.4–7.8 AU 之间、寿命较短的“轨道通道”，约有 21\% 的半人马天体会经过该区域，其中包括 72\% 最终演化为木星族彗星的半人马天体。目前已知有 4 个天体位于这一通道区域内（29P/Schwassmann–Wachmann、P/2010 TO20 LINEAR-Grauer、P/2008 CL94 Lemmon 以及 2016 LN8），但数值模拟表明，可能还存在数量级约为 1000 个、半径大于 1 km、尚未被探测到的天体。处于该通道区域内的天体可能表现出显著的活动性$^{[17][18]}$，并且正处于一个重要的演化过渡状态，这进一步模糊了半人马天体与木星族彗星两类天体之间的界限。
\subsection{命名惯例}
根据国际天文学联合会小天体命名工作组（Working Group for Small Bodies Nomenclature，WGSBN；原小天体命名委员会）$^{[19]}$ 的规定，半长轴小于 30 AU（位于海王星轨道以内）且近日点距离大于 5.5 AU（位于木星轨道之外）的半人马天体，应以神话中的半人马来命名，即半人半马的神话生物$^{[20]}$。WGSBN 在技术上将穿越海王星轨道的跨海王星天体（半长轴大于 30 AU、近日点小于 30 AU）也计入半人马天体的范畴，但为其保留了另一套不同的命名方案$^{[20]}$。该方案于 2007 年被正式采用，当时首批此类天体——65489 Ceto–Phorcys 和 42355 Typhon–Echidna——被命名$^{[19]}$。

根据 WGSBN 的规定，穿越海王星轨道的跨海王星天体必须以神话中的奇美拉类生物命名$^{[20]}$，这类生物包括具有混合形态或可变形态的神话生物$^{[19]}$。这一类别中的一个命名实例是 471325 Taowu，其名称源自中国神话中的“饕餮”，据说是一种兼具人、虎和野猪特征的混合生物$^{[21]}$。
\subsection{轨道}
\subsubsection{分布}
\begin{figure}[ht]
\centering
\includegraphics[width=6cm]{./figures/acc469512e466fac.png}
\caption{已知半人马小行星的轨道分布$^{[note 2]}$} \label{fig_Centau_3}
\end{figure}
该示意图展示了已知半人马天体的轨道相对于行星轨道的位置关系。对于部分选定的天体，其轨道偏心率以红色线段表示（从近日点延伸至远日点）。

半人马天体的轨道偏心率分布范围很广，从高度偏心的轨道（如 Pholus、Asbolus、Amycus、Nessus）到接近圆形的轨道（如 Chariklo 以及穿越土星轨道的 Thereus 和 Okyrhoe）。

为说明轨道参数的多样性，图中以黄色标出了若干具有极端轨道特征的天体：
\begin{itemize}
\item 1999 XS35（阿波罗型小行星）具有极端的高偏心率轨道（(e = 0.947)），其轨道范围从地球轨道以内（0.94 AU）一直延伸到远超海王星轨道之外（$>34 AU$）。
\item 2007 TB434 具有近乎圆形的轨道（$(e < 0.026)$）。
\item 2001 XZ255 具有最低的轨道倾角（$(i < 3^\circ)$）。
\item 2004 YH32 是少数具有极端顺行高倾角（$(i > 60^\circ)$）的半人马天体之一，其轨道倾角高达 $(79^\circ)$。尽管其轨道在太阳距离上从小行星带附近延伸至土星轨道以外，但若将其轨道投影到木星轨道平面上，甚至不会到达木星的距离。
\end{itemize}
已有十余个已知半人马天体沿逆行轨道运行，其轨道倾角从中等值（例如 Dioretsa 的约 $(160^\circ)$）,到极端值（$(i>120^\circ)$，例如 (342842) 2008 YB3 的约,$(105^\circ)$ $^{[22]}$）。其中有十七个高倾角逆行半人马天体曾被有争议地认为可能具有星际起源 $^{[23][24][25]}$。
\subsubsection{轨道变化}
\begin{figure}[ht]
\centering
\includegraphics[width=8cm]{./figures/8ef3021d0bf5acc7.png}
\caption{在未来 5500 年内，Asbolus 的半长轴变化情况，采用了两组略有不同的当前轨道要素估计进行计算。在公元 4713 年与木星发生近距离遭遇之后，这两种计算结果开始出现分歧。$^{[26]}$} \label{fig_Centau_4}
\end{figure}
由于半人马天体不受轨道共振的保护，它们的轨道在 $10^{6}-10^{7}$ 年的时间尺度内是不稳定的。$^{[27]}$ 例如，55576 Amycus 位于接近天王星 3:4 共振的位置，其轨道是不稳定的。$^{[1]}$ 对其轨道的动力学研究表明，处于“半人马天体”状态很可能是天体从柯伊伯带向木星族短周期彗星过渡过程中的一个中间轨道阶段。小行星 ((679997)) 2023 RB 将在 2201 年与土星的一次近距离接近中，其轨道发生显著改变。

天体可能首先从柯伊伯带受到扰动，从而成为穿越海王星轨道的天体，并与该行星发生引力相互作用（见其起源理论）。随后它们被归类为半人马天体，但其轨道呈现混沌特征，在与一个或多个外行星反复发生近距离接近的过程中迅速演化。一些半人马天体会进一步演化为穿越木星轨道的天体，其近日点可能被降低至内太阳系；若此时它们表现出彗星活动，则可能被重新归类为木星族活动彗星。最终，半人马天体要么与太阳或某一行星发生碰撞，要么在与行星（尤其是木星）的一次近距离接近后被抛射进入星际空间。
\subsection{物理特性}
\begin{figure}[ht]
\centering
\includegraphics[width=6cm]{./figures/7822bdb4c4d1ac8a.png}
\caption{土星的卫星菲比（Phoebe）于 2004 年由卡西尼–惠更斯号探测器成像，理论上被认为是一颗被俘获的半人马天体。} \label{fig_Centau_5}
\end{figure}
与矮行星和小行星相比，半人马天体的尺寸相对较小且距离较远，这使得对其表面的遥感观测变得不可行；然而，颜色指数和光谱仍可为其表面成分提供线索，并有助于理解这些天体的起源。$^{[27]}$
\subsubsection{颜色}
\begin{figure}[ht]
\centering
\includegraphics[width=6cm]{./figures/2848d3a0f02059ff.png}
\caption{半人马天体的颜色分布} \label{fig_Centau_6}
\end{figure}
半人马天体的颜色具有高度多样性，这对任何关于其表面成分的简单模型都构成了挑战。$^{[28]}$ 在旁图中，颜色指数是通过蓝光（B）、可见光（V）（即绿—黄色）和红光（R）滤光片测得的视星等差值。该图以夸张的颜色展示了所有已知颜色指数的半人马天体之间的差异。作为参照，还绘制了两颗卫星——海卫一（Triton）和菲比（Phoebe），以及行星火星（黄色标注，尺寸未按比例）。

半人马天体似乎可以分为两类：
\begin{itemize}
\item 一类是非常偏红的，例如 5145 Pholus；
\item 另一类是偏蓝的（或按部分作者的说法为蓝灰色），例如 2060 Chiron 或 2020 MK4。
\end{itemize}
关于这种颜色差异的解释存在多种理论，但大体可以归为两类：
\begin{itemize}
\item 其一，颜色差异源于半人马天体在起源或组成上的不同（见下文“起源”）；
\item 其二，颜色差异反映了它们受到辐射和/或彗星活动影响的空间风化程度不同。
\end{itemize}
作为第二类解释的例子，Pholus 的红色被解释为其表面可能覆盖着一层经辐照的红色有机物幔层；而 Chiron 由于其周期性的彗星活动，使冰层暴露在表面，从而呈现出蓝色或灰色的颜色指数。然而，活动性与颜色之间的相关性并不确定，因为已知具有活动性的半人马天体颜色范围很广，从蓝色的 Chiron 到红色的 166P/NEAT 均有分布。$^{[29]}$ 另一种可能是，Pholus 只是最近才从柯伊伯带被抛射出来，因此其表面转化过程尚未发生。

Delsanti 等人提出，可能同时存在多种相互竞争的过程，例如辐射导致的表面变红，以及碰撞引起的“变浅”效应。$^{[30][31]}$
\subsubsection{光谱}
光谱的解释往往具有不确定性，这与颗粒尺寸及其他因素有关，但光谱仍然能够为表面成分提供重要线索。与颜色分析类似，观测到的光谱往往可以符合多种不同的表面模型。

在若干半人马天体上已经确认存在水冰的光谱特征$^{[27]}$，其中包括 2060 Chiron、10199 Chariklo 和 5145 Pholus。除水冰特征之外，还提出了多种其他表面组成模型：

有人认为 Chariklo 的表面是索林（tholins，类似于在土卫六和海卫一上探测到的物质）与无定形碳的混合物。
Pholus 的表面被认为覆盖着一种类似土卫六的索林、炭黑、橄榄石$^{[32]}$以及甲醇冰的混合物。
52872 Okyrhoe 的表面被认为由干酪根、橄榄石以及少量水冰组成。
8405 Asbolus 的表面可能由 15\% 类似海卫一的索林、8\% 类似土卫六的索林、37\% 无定形碳以及 40\% 冰态索林构成。
Chiron 的情况似乎最为复杂，其观测到的光谱会随观测时期而变化。在低活动期探测到了水冰特征，而在高活动期这一特征则会消失。$^{[33][34][35]}$
\subsubsection{与彗星的相似性}
\begin{figure}[ht]
\centering
\includegraphics[width=6cm]{./figures/7eeab0f5e3ea6443.png}
\caption{彗星 38P 在 1982 年至 2067 年间多次与木星、土星和天王星发生近距离接近，表现出类似半人马天体的行为。$^{[36]}$} \label{fig_Centau_7}
\end{figure}
对 Chiron 在 1988 年和 1989 年接近其近日点时的观测发现，它表现出了彗发（即从其表面蒸发出的气体和尘埃云）。因此，Chiron 现已被正式同时归类为小行星和彗星，尽管它的尺寸远大于典型彗星，这一点仍存在一定争议。其他半人马天体也在被持续监测其是否存在类似彗星的活动：迄今为止，已有两颗——60558 Echeclus 和 166P/NEAT——表现出这种行为。166P/NEAT 在被发现时就已呈现彗发，因此被归类为彗星，尽管其轨道属于半人马天体。Echeclus 被发现时并无彗发，但近期开始表现出活动性$^{[37]}$，因此如今也被同时归类为彗星和小行星。总体而言，大约有 30 颗半人马天体被探测到存在活动性，而这一活动性天体的样本明显偏向近日点距离较小的对象$^{[38]}$。

在 Echeclus$^{[8]}$ 和 Chiron$^{[39]}$ 中均探测到了极少量的一氧化碳，其推算出的一氧化碳释放率足以解释所观测到的彗发。然而，这两者的一氧化碳释放率显著低于通常在 29P/Schwassmann–Wachmann$^{[17]}$ 中观测到的水平，后者是一颗远距离仍保持活动的彗星，常被归类为半人马天体。

在轨道特性上，半人马天体与彗星之间并不存在明确的区分。29P/Schwassmann–Wachmann 和 39P/Oterma 都曾被称为半人马天体，因为它们具有典型的半人马轨道。彗星 39P/Oterma 目前处于非活动状态，其仅在 1963 年被木星扰动进入半人马轨道之前表现出活动性$^{[40]}$。暗弱的彗星 38P/Stephan–Oterma 如果其近日点距离位于木星轨道之外、约 5 AU 处，很可能也不会表现出彗发。到 2200 年左右，彗星 78P/Gehrels 可能会向外迁移，进入类似半人马天体的轨道
\subsubsection{自转周期}
对 Chiron 和 Chariklo 的光变曲线进行周期图分析，分别得到它们的自转周期为$5.5 \pm 0.4$小时和$7.0 \pm 0.6$小时。$^{[41]}$
\subsubsection{尺寸、密度与反照率}
\begin{figure}[ht]
\centering
\includegraphics[width=8cm]{./figures/9608c0bd2bfccf0e.png}
\caption{具有已测量直径的多颗大型半人马天体在尺寸、反照率和颜色方面的对比图。} \label{fig_Centau_8}
\end{figure}
半人马天体的直径可达数百千米。最大的半人马天体直径超过 300 千米，并且主要分布在 20 AU 以外的区域。$^{[42]}$
\subsection{起源假说}
关于半人马天体起源的研究近年来取得了许多新进展，但由于物理观测数据仍然有限，任何结论都受到一定限制。目前已提出多种可能的起源模型。

数值模拟表明，一些柯伊伯带天体的轨道可能受到扰动，从而被抛射出来并演化为半人马天体。散射盘天体在动力学上是这种抛射过程的最佳候选者（例如，半人马天体可能属于从柯伊伯带向内扰动而来的“内侧”散射盘天体）。然而，散射盘天体的颜色并不符合半人马天体所呈现出的双色特征。冥族天体（Plutinos）是柯伊伯带天体中的一类，表现出类似的双色特性，因此有人提出，并非所有冥族天体的轨道都如最初认为的那样稳定，这可能与冥王星的引力扰动有关。$^{[43]}$ 随着对柯伊伯带天体获得更多物理数据，预计这一问题将有进一步进展。

部分半人马天体可能起源于碎裂事件，这些事件或许是在与木星的近距离接近过程中被触发的。$^{[44]}$ 半人马天体 2020 MK4、P/2008 CL94（Lemmon）和 P/2010 TO20（LINEAR-Grauer）的轨道都与彗星 29P/Schwassmann–Wachmann 的轨道非常接近。29P 是最早被发现的半人马天体，在其处于活动状态时，这些天体之间可能发生近距离接近，甚至其中某个天体会穿过 29P 的彗发。$^{[44]}$

至少有一颗半人马天体，即 2013 VZ70，可能起源于土星不规则卫星族群，其形成机制可能涉及撞击、碎裂或潮汐破坏。$^{[45]}$
\subsection{著名的半人马天体}
\begin{figure}[ht]
\centering
\includegraphics[width=10cm]{./figures/96dcaeb9918758ed.png}
\caption{} \label{fig_Centau_9}
\end{figure}
a. 该类别由物体的近日点和远日点距离定义: S表示土星附近的近日点/远日点，天王星附近的U，海王星附近的N和柯伊伯带中的K。
\subsection{另见}
\begin{itemize}
\item 小行星
\item 矮行星
\end{itemize}
\subsection{说明性注释}
\begin{enumerate}
\item 有关此想法的批评$^{[3]}$
\item 为了这个图的目的，一个对象被归类为半人马，如果它的半长轴位于木星和海王星之间
\end{enumerate}
\subsection{参考资料}
\begin{enumerate}
\item 霍纳，J.；西北埃文斯；Bailey, M.E.(2004)。“半人马种群的模拟I: 批量统计”。皇家天文学会月刊。354(3):798-810。arXiv:astro-ph/0407400。Bibcode:2004MNRAS.354..798H。doi:10.1111/j.1365-2966.2004.08240.x。S2CID 16002759。
\item Fathi Namouni和Maria Helena Moreira Morais (2018年5月2日)。“木星逆行同轨道小行星的星际起源”。皇家天文学会月刊。477(1):L117-L121。arXiv:1805.09013。Bibcode:2018MNRAS.477L.117N。doi:10.1093/mnrasl/sly057。S2CID54224209。  
\item Billings，Lee (2018年5月21日)。天文学家发现潜在的 “星际” 小行星围绕太阳向后运行。科学美国人。已检索6月1日2018。
\item 萨里德，G.；沃尔克，K.；斯特克洛夫，J.；哈里斯，W.；沃马克，M.；伍德尼，L.(2019年)。“29P/schwassmann-wachmann 1，通往木星家族彗星的半人马”。天体物理学杂志快报。883(1): 7。arXiv:1908.04185。Bibcode:2019ApJ...883L..25S。doi:10.3847/2041-8213/ab3fb3。S2CID199543466。 
\item 谢泼德，S.；Jewitt, D.;特鲁希略，C.；布朗，M.；Ashley, M.(2000)。“半人马和柯伊伯带天体的广域CCD调查”。天文杂志。120(5):2687-2694。arXiv:astro-ph/0008445。Bibcode:2000AJ....120.2687S。doi:10.1086/316805。S2CID119337442。 
\item David Jewitt; Nader Haghighipour (2007)。“行星的不规则卫星: 早期太阳系捕获的产物”(PDF)。天文学和天体物理学年度评论。45(1):261-95。arXiv:astro-ph/0703059。Bibcode:2007ARA & A .. 45 .. 261J。doi:10.1146/annurev.astro.44.051905.092459。S2CID13282788。存档自原来的(PDF)在2009-09-19。  
\item “谷神星的黎明: 我们学到了什么？”(PDF)。空间研究委员会。国家科学院。存档自原来的(PDF)关于2018-04-13。已检索2023-10-11。
\item Wierzchos, K.;沃马克，M.；Sarid, G.(2017)。“远距离活跃的半人马 (60558) 174p/Echeclus中的一氧化碳，位于6 au”。天文杂志。153(5): 8。arXiv:1703.07660。Bibcode:2017AJ....153..230W。doi:10.3847/1538-3881/aa689c。S2CID119093318。 
\item “不寻常的小行星”。小行星中心。已检索10月25日2010。
\item “半人马和分散磁盘对象列表”。小行星中心。已检索2025-10-01。
\item “轨道分类 (半人马座)”。JPL太阳能系统动力学。存档自原来的2021年8月1日。已检索10月13日2008。
\item 艾略特，J.L.；科恩，S.D.;克兰西，K。B.;古尔比斯，A。A.美国；米利斯，R。L.;Buie, M.W.;Wasserman, L.H、; 蒋，E。I.;约旦，A。B.;特里林，D。E.;Meech, K.J.(2005)。“深黄道调查: 寻找柯伊伯带天体和半人马。二。动态分类，柯伊伯带平面和核心种群“。天文杂志。129(2):1117-1162。Bibcode:2005AJ....129.1117E。doi:10.1086/427395。
\item 格拉德曼，B。;马斯登，B。范·拉尔霍芬，C。(2008)。外太阳系 (海王星以外的太阳系) 的命名(PDF)。亚利桑那大学出版社。ISBN  978-0-8165-2755-7。存档自原来的 (PDF)在2012-11-02。
\item 蔡恩，尤金; 利维克，Y.；默里-克莱，R.；Buie, M.;格伦迪，W.；霍尔曼，M。(2007)。Reipurth, B.;Jewitt, D.;凯尔，K. (编辑)。“跨海王星空间简史”。原恒星和行星V。亚利桑那大学出版社:895-911.arXiv:astro-ph/0601654。Bibcode:2007prpl.conf .. 895C。doi:10.1086/316805。
\item “JPL小型数据库搜索引擎: 半人马列表”。JPL太阳能系统动力学。存档自原来的2021年4月20日。已检索10月11日2018。
\item “JPL小体数据库搜索引擎: 比天王星的轨道更近的TNOs列表”。JPL太阳能系统动力学。存档自原来的2021年8月29日。已检索10月11日2018。
\item 沃马克，M.；Wierzchos, K.;Sarid, G.(2017)。“遥远活跃的彗星中的CO”。太平洋天文学会的出版物。129(973) 031001。arXiv:1611.00051。Bibcode:2017PASP .. 129c1001W。doi:10.1088/1538-3873/129/973/031001。S2CID118507805。 
\item Lacerda，P。(2013)。“彗星P/2010 TO20线性-格劳尔作为Mini-29P/SW1”。皇家天文学会月刊。883(2):1818-1826年。arXiv:1208.0598。Bibcode:2013MNRAS.428.1818L。doi:10.1093/mnras/sts164。S2CID54030926。 
\item 格伦迪，威尔; 斯坦斯伯里，J.A.；诺尔，K; 斯蒂芬斯，华盛顿特区；特里林，D.E.；克恩，s.d.；斯宾塞，J.R.；克鲁克申克，d.p.；利维森，h.f.(2007)。“(65489) Ceto/Phorcys的轨道，质量，大小，反照率和密度: 潮汐进化的二进制半人马座”。伊卡洛斯。191(1):286-297。arXiv:0704.1523。Bibcode:2007年Icar。。191。。286克。doi:10.1016/j.icarus.2007.04.004。S2CID1532765。 
\item “命名非彗星小型太阳系天体的规则和准则”(PDF)。IAU小型机构命名工作组。2025年2月22日。已检索2025-09-10。
\item “JPL小体数据库查找: 471325 Taowu (2011 KT19)”(2024-05-27最后一项。)。喷气推进实验室。已检索2025-09-13。
\item C.德拉富恩特·马科斯; R.德拉富恩特·马科斯 (2014年)。“大型逆行半人马: 来自奥尔特云的游客？”天体物理学和空间科学。352(2):409-419。arXiv:1406.1450。Bibcode:2014Ap & SS.352..409D。doi:10.1007/s10509-014-1993-9。S2CID119255885。 
\item Fathi Namouni和Maria Helena Moreira Morais (2020年5月)。“高倾角半人马的星际起源”。皇家天文学会月刊。494(2):2191-2199。arXiv:2004.10510。Bibcode:2020MNRAS.494.2191N。doi:10.1093/mnras/staa712。S2CID216056648。 
\item Raymond, S.N.;布拉塞尔，R.；巴蒂金，K.；莫比德利，A.(2020)。“没有证据表明星际小行星被困在太阳系中”。皇家天文学会每月通知: 信件。497(1):L46-L49。arXiv:2006.04534。Bibcode:2020MNRAS.497L..46M。doi:10.1093/mnrasl/slaa111。S2CID219531537。  
\item 纳穆尼，法蒂 (2022年)。“行星穿越小行星的倾斜路径”。皇家天文学会月刊。510:276-291。arXiv:2111.10777。doi:10.1093/mnras/stab3405。
\item “centaur 8405 Asbolus的三个克隆在450Gm内通过”。存档自原来的on 2015-09-13。已检索2009-05-02。("Solex 10"。存档自原来的在2008-12-20。)
\item Jewitt, David C.; A.德尔桑蒂 (2006年)。“超越行星的太阳系”。太阳系更新: 太阳系科学的专题和及时评论。施普林格-实践版。ISBN 978-3-540-26056-1。(预印本 (pdf))
\item 巴鲁奇，M.A.;Doressoundiram, A.;克鲁克申克，D。P。(2003)。“TNOs和半人马的物理特征”(PDF)。巴黎天文台空间研究和天体物理学仪器实验室。存档自原来的(PDF)2008年5月29日。已检索3月20日2008。
\item 詹姆斯·M·鲍尔；Fernández, Yanga R.;Meech, Karen J.(2003)。“活跃的半人马C/NEAT的光学调查 (2001 T4)”。太平洋天文学会的出版物。115(810):981-989。Bibcode:2003PASP..115..981B。doi:10.1086/377012。S2CID122502310。 
\item Peixinho, N.;Doressoundiram, A.;德尔桑蒂，A.；博恩哈特，H.; 巴鲁奇，M。A.;Belskaya，我.(2003)。“重新开启TNOs颜色争议: 半人马双峰和TNOs单峰”。天文学和天体物理学。410(3):L29-L32。arXiv:astro-ph/0309428。Bibcode:2003A & A...410L..29P。doi:10.1051/0004-6361:20031420。S2CID8515984。  
\item Hainaut & Delsanti (2002)外太阳系小天体的颜色天文学和天体物理学，389,641数据源
\item 镁铁硅酸盐 (Mg，Fe)2SiO4,常见的组件岩岩石。
\item Dotto，E; Barucci，M A; De Bergh，C (2003年6月)。“半人马的颜色和组成”。地球、月球和行星。92(1-4):157-167。Bibcode:2003EM & P...92 .. 157D。doi:10.1023/b: 月0000031934.89097.88。S2CID189905595。  
\item Luu, Jane X.;Jewitt, David;特鲁希略，C。A.(2000)。“2060年凯龙星的水冰及其对半人马和柯伊伯带天体的影响”。天体物理学杂志。531(2):L151-L154。arXiv:astro-ph/0002094。Bibcode:2000ApJ...531L.151L。doi:10.1086/312536。PMID10688775。S2CID9946112。   
\item 费尔南德斯，Y.R、；Jewitt, D.C.；谢泼德，S。S。(2002)。“半人马Asbolus和Chiron的热特性”。天文杂志。123(2):1050-1055。arXiv:astro-ph/0111395。Bibcode:2002AJ....123.1050F。doi:10.1086/338436。S2CID11266670。 
\item “JPL接近数据: 38p/斯蒂芬-奥特马”。NASA.1981-04-04.最后一个obs。存档自原来的on 2021-09-21。已检索2009-05-07。
\item 崔，y-j.；魏斯曼，P.R.；波利沙克，D。(2006年1月)。“(60558) 2000 EC_98”。IAU Circ。(8656): 2。Bibcode:2006IAUC.8656....2C。
\item Jewitt, D.(2009)。“活跃的半人马”。天文杂志。137(5):4295-4312。arXiv:0902.4687。Bibcode:2009AJ....137.4296J。doi:10.3847/1538-3881/aa689c。S2CID119093318。 
\item 沃马克，M.；斯特恩，A.(1999)。“(2060年) Chiron中一氧化碳的观察”(PDF)。月球与行星科学二十八。已检索2017-07-11。
\item Mazzotta Epifani, E.;帕伦博，P.；卡普里亚，M。T.;克雷蒙人，G.；富勒，M.；Colangeli, L.(2006)。“活跃的半人马P/2004 A1 (LONEOS) 的尘埃昏迷: 共同驱动的环境？”。天文学与天体物理学。460(3):935-944。Bibcode:2006A & A...460 .. 935米。doi:10.1051/0004-6361:20065189。
\item Galiazzo, M.A.; 德拉富恩特·马科斯，C.; 德拉富恩特·马科斯，R.；卡拉罗，G.；马里斯，M.；蒙塔尔托，M。(2016)。“半人马和跨海王星物体的光度法: 2060 Chiron (1977 UB)，10199 Chariklo (1997 CU26)，38628 Huya (2000 EB173)，28978 Ixion (2001 KX76)，和90482奥库斯 (2004 DW)“。天体物理学和空间科学。361(3):212-218.arXiv:1605.08251。Bibcode:2016Ap & SS.361..212G。doi:10.1007/s10509-016-2801-5。S2CID119204060。 
\item Galiazzo, M.A.;维格特，P。&Aljbaae, S.(2016)。“半人马和TNOs对主带及其家族的影响”。天体物理学和空间科学。361(12):361-371。arXiv:1611.05731。Bibcode:2016Ap & SS.361..371G。doi:10.1007/s10509-016-2957-z。S2CID118898917。 
\item 万，X.-S。；黄，T.-Y。(2001)。“超过1亿年的32个plutinos的轨道演化”。天文学和天体物理学。368(2):700-705。Bibcode:2001A & A...368 .. 700瓦。doi:10.1051/0004-6361:20010056。
\item de la Fuente Marcos, C.; de la Fuente Marcos, R.;利桑德罗，J.；Serra-Ricart, M.;马蒂诺，S.; 德莱昂，J.；Chaudry, F.;阿拉尔孔，M。R。(2021年5月13日)。“活跃的半人马2020 MK4”。天文学与天体物理学。649(1): A85 (15页)。arXiv:2104.01668。Bibcode:2021A & A...649A..85D。doi:10.1051/0004-6361/202039117。S2CID233024896。 
\item de la Fuente Marcos, C.; de la Fuente Marcos, R.(2022年1月10日)。“半人马2013 VZ70: 来自土星不规则月球种群的碎片？”。天文学与天体物理学。657(1): A59 (10页)。arXiv:2110.04264。Bibcode:2022A & A...657A..59D。doi:10.1051/0004-6361/202142166。S2CID238856647。 
 \end{enumerate}